\newpage
\phantomsection
\label{prf:differentiable-implies-continuous}

\begin{remark*}[Return]
\hyperref[thm:differentiable-implies-continuous]{Return to Theorem}
\end{remark*}

\begin{theorem*}[Differentiability Implies Continuity]
Let $f$ be defined on an open interval $I$ containing $c$. If $f$ is differentiable at $c$, then $f$ is continuous at $c$.
\end{theorem*}

\begin{proof}
\textbf{Professional Standard Proof.}~\\
Assume $f$ is differentiable at $c$. For any $x \in I$ with $x \neq c$, the identity $f(x) - f(c) = \frac{f(x) - f(c)}{x - c} \cdot (x - c)$ holds. Taking limits as $x \to c$ and applying the Product Rule for Limits yields $\lim_{x \to c}[f(x) - f(c)] = \lim_{x \to c}\frac{f(x) - f(c)}{x - c} \cdot \lim_{x \to c}(x - c) = f'(c) \cdot 0 = 0$. Therefore $\lim_{x \to c} f(x) = f(c)$, establishing continuity at $c$.
\end{proof}

\begin{proof}
\textbf{Detailed Learning Proof.}~\\

\textbf{Step 1.} Factor the vertical displacement using secant identity. For any $x \in I$ with $x \neq c$, the vertical displacement admits the factorization $f(x) - f(c) = \left(\frac{f(x) - f(c)}{x - c}\right) \cdot (x - c)$. This identity holds on all of $I \setminus \{c\}$ by arithmetic; no limit has been taken.

\textbf{Step 2.} Distribute the limit over the product. Taking the limit as $x \to c$ and applying the Product Rule for Limits: $\lim_{x \to c} [f(x) - f(c)] = \left(\lim_{x \to c} \frac{f(x) - f(c)}{x - c}\right) \cdot \left(\lim_{x \to c} (x - c)\right)$. The Product Rule applies because both limits on the right exist --- the first by the differentiability hypothesis, the second trivially.

\textbf{Step 3.} Evaluate each factor. By definition of differentiability at a point, the first factor equals $f'(c)$, a finite real number. The second factor satisfies $\lim_{x \to c}(x - c) = 0$ by the Identity Limit. Therefore: $\lim_{x \to c} [f(x) - f(c)] = f'(c) \cdot 0 = 0$.

\textbf{Step 4.} Conclude continuity. Since $\lim_{x \to c}[f(x) - f(c)] = 0$, it follows that $\lim_{x \to c} f(x) = f(c)$. By definition of continuity at a point, $f$ is continuous at $c$.
\end{proof}

\begin{remark*}[Proof structure]
The proof employs a direct approach, hinging on the algebraic factorization of $f(x) - f(c)$ into the product of the difference quotient and the displacement. This converts the continuity problem into a limit-of-product problem, where the derivative hypothesis controls one factor and the displacement annihilates the product.
\end{remark*}

\begin{remark*}[Dependencies]~\\
\begin{itemize}
  \item \hyperref[def:derivative]{Definition of Differentiability at a Point}
  \item \hyperref[prop:limit-product-rule]{Product Rule for Limits}
  \item \hyperref[prop:identity-limit]{Identity Limit}
  \item \hyperref[def:continuity]{Definition of Continuity at a Point}
\end{itemize}
\end{remark*}
